\documentclass[11pt]{article}
\usepackage[margin=1in]{geometry}
\usepackage{amsmath,amssymb,amsthm,booktabs}
\usepackage[hidelinks]{hyperref}
\setlength{\emergencystretch}{2em}
\newtheorem{theorem}{Theorem}
\newtheorem{proposition}[theorem]{Proposition}
\newtheorem{corollary}[theorem]{Corollary}
\newcommand{\mr}{\operatorname{mr}}
\newcommand{\Sym}{\operatorname{Sym}}
\newcommand{\GL}{\operatorname{GL}}
\newcommand{\Z}{\mathbb Z}
\newcommand{\C}{\mathbb C}
\newcommand{\F}{\mathbb F}
\title{Six-point covers and formal orthogonality at icosahedral points}
\author{Internal research working proof}
\date{16 November 2026}
\begin{document}
\maketitle

The six-point action of $A_5$ gives a different augmentation module
from its five-point action and permits characteristic five in the
split homology-gain test. All sixteen retained $A_5$ quotient-kernel
classes on the current sixteen-knot panel have now been checked in
this action. Their index-six subgroup homology is cyclic of order
$4D$ or $28D$. The maximum gain is one, short of the required six.

These finite homology computations also supply a stronger structural
boundary. For the associated three-dimensional icosahedral reflection
representations of the knot orbifold group, \emph{every formal
characteristic-zero deformation preserving the meridian-square
relation continues to preserve a nondegenerate symmetric form}.
The longitude eigenvalue on the meridian's nonidentity line therefore
still satisfies $\lambda^2=1$. Thus nonlinear deformation through
these points with meridians kept involutive cannot escape the
peripheral obstruction of the companion augmentation note. One must
allow the meridian eigenvalue to vary, use another base component,
or change representation dimension. This does not exclude those
possibilities or settle a new knot. Full MRC remains open; the
uniform proofs below are working proofs without formal verification,
external specialist review or a novelty claim.

\section{The six-point action and exact cover homology}

Let $G=\pi_1(\Sigma_2K)$ and $\phi:G\twoheadrightarrow A_5$ be one
of the retained quotient maps. Conjugation on the six Sylow-five
subgroups gives a faithful transitive action of degree six. The
point stabilizer is a normalizer $D_5$ of order ten; set
$H_6=\phi^{-1}(D_5)$. All six such subgroups are conjugate, and
automorphisms of $A_5$ permute them. One epimorphism per quotient
kernel therefore suffices for these unbased covers.

The complex augmentation character $\chi_5$ has values
$5,1,-1,0,0$ on the classes $1,2,3,5A,5B$ of sizes
$1,15,20,12,12$. Its norm is $(25+15+20)/60=1$, so it is
irreducible. The old degree-five augmentation character has values
$4,0,1,-1,-1$, with inner product zero against $\chi_5$.
This is a new coefficient module, rather than a copy of the old one.

For every prime $p\nmid6$, the general augmentation bound in the
companion note gives
\[
 b_1(H_6;\F_p)-b_1(G;\F_p)\ge6
 \quad\Longrightarrow\quad \operatorname{rank}G\ge3
 \quad\Longrightarrow\quad \mr(K)\ge4.
\]
The new panel was frozen before computing. It uses all sixteen
kernel classes on exactly the old sixteen targets, with no adaptive
map or target replacement. Every base has $H_1(G;\Z)=\Z/D$.
The table lists the subgroup homology for \emph{each} knot in a row.
When two groups are listed, they belong to its two separate kernels.
\begin{center}
\begin{tabular}{lrl}
\toprule
Knot or knots & $D$ & Index-six subgroup homology\\
\midrule
$13a129,13a136,13a194,13a197$&391&$\Z/(4D)$\\
$13a135,13a137,13a196,13a198$&401&$\Z/(4D)$ and $\Z/(28D)$\\
$13a234$&389&$\Z/(28D)$\\
$13a242$&359&$\Z/(28D)$\\
$13a244,13a265$&311&$\Z/(28D)$\\
\bottomrule
\end{tabular}
\end{center}
The remaining four targets have no $A_5$ quotient. All listed $D$
are coprime to $210$. The $4D$ cases give gain zero for $p\nmid6$;
the $28D$ cases give gain one at seven and zero at every other
allowed prime. In particular characteristic five gives no gain.
These statements follow from integral Smith forms, not a prime sweep.
There are eight cases of each type, on separate source groups or
kernels; their maps are never pooled into a joint quotient count.

This exhausts the $A_5$-induced degree-six action on these targets.
It does \emph{not} exhaust every degree-six cover: other monodromy,
including $A_6$, is outside this panel.

\section{The exceptional characteristics and knot descent}

Write $V_{n-1}=\ker(F^n\to F)$ for the augmentation module. If
$p\mid n$, its constant vector is invariant and the splitting used
by the Betti-gain formula fails. The beginning of the long exact
cohomology sequence of
\[
             0\longrightarrow V_{n-1}\longrightarrow F^n
               \longrightarrow F\longrightarrow0
\]
handles the current targets anyway. Transitivity gives one-dimensional
invariants in both $V_{n-1}$ and $F^n$, and the sum map on constants
is zero. If $H^1(G,F)=0$, exactness and Shapiro give
\[
 \dim H^1(G,V_{n-1})=1+b_1(H_n;F).
\]
Here the base Betti number vanishes at $2,3,5$. Thus the previously
omitted characteristic five for $V_4$ and characteristics two and
three for $V_5$ have the following exact dimensions:
\begin{center}
\begin{tabular}{lrrrr}
\toprule
Module and prime & $\dim V^G$ & $\dim B^1$ & $\dim H^1$ & $\dim Z^1$\\
\midrule
$V_4$, $p=5$ &1&3&1&4\\
$V_5$, $p=2$ &1&4&2&6\\
$V_5$, $p=3$ &1&4&1&5\\
\bottomrule
\end{tabular}
\end{center}
Direct cocycle matrices independently reproduce these dimensions for
every class. Combining them with the split cases shows that neither
the $V_4$ nor the $V_5$ ordinary-generation cocycle estimate can force
$\operatorname{rank}G\ge3$ in \emph{any} coefficient characteristic.
This is a boundary for these modules, not a proof that the groups
are two-generated or a statement about every $A_5$ module.

Let $Q=G_K/\langle\!\langle\mu^2\rangle\!\rangle$ be the knot
orbifold group. It contains $G$ normally with index two. Every tested
map has a unique involution $t\in A_5$ implementing the deck
automorphism on its arc generators: $t\phi(g_i)t=\phi(g_i)^{-1}$.
Consequently the meridian images $\phi(g_i)t$ define the six-point
coefficient representation of $Q$. These assertions are checked
directly in the separate matrix model.

For $p\ge5$, the direct orbifold cocycle matrix gives
$H^1(Q,V_5)=0$ for all sixteen classes. The only possibly nonzero
cover contribution is at seven; the eight cover classes there do
not descend. At all other primes $p\ge5$, restriction into the
already-zero cover cohomology proves vanishing. Every admissible
affine lift on this fixed module is therefore principal. This
argument does not impose the old three-dimensional root geometry;
it computes the new module separately. We do not assert this
orbifold statement in characteristics two or three.

\section{A uniform lifting lemma for invariant forms}

The following standard cohomological mechanism supplies a nonlinear
consequence of the degree-six homology computation.

\begin{theorem}\label{thm:forms}
Let $\rho_0:Q\to\GL(W)$ over a field $k$ of characteristic different
from two preserve a nondegenerate
symmetric form $B_0$, and put $S=\Sym^2(W^*)$ with action
\[
                    g\cdot B=\rho_0(g)^{-\mathsf T}B\rho_0(g)^{-1}.
\]
If $H^1(Q,S)=0$, every formal representation
$\rho_t:Q\to\GL(W\otimes_k k[[t]])$ reducing to $\rho_0$
preserves a symmetric form $B(t)$ with $B(0)=B_0$.
The resulting form is nondegenerate over $k[[t]]$.
\end{theorem}
\begin{proof}
Suppose a symmetric form $B^{(j)}$ with constant term $B_0$ is
invariant modulo $t^j$; start with $B^{(1)}=B_0$. Let $A_t(g)$ be
the action on symmetric forms induced by $\rho_t$. The leading
failure of invariance is
\[
              c_j(g)=t^{-j}\bigl(A_t(g)B^{(j)}-B^{(j)}\bigr)\bmod t.
\]
The identity $A_t(gh)=A_t(g)A_t(h)$ shows
$c_j(gh)=c_j(g)+g\cdot c_j(h)$, so $c_j$ is an $S$-valued
cocycle. By the hypothesis it equals $g\cdot C_j-C_j$ for some
$C_j\in S$. Replacing $B^{(j)}$ by
$B^{(j+1)}=B^{(j)}-t^jC_j$ kills that failure modulo $t^{j+1}$.
These successive corrections converge in $k[[t]]$ to an invariant
symmetric form. Its determinant has the nonzero constant term
$\det B_0$, so it is a unit.
\end{proof}

This proof uses cohomology of the original module at $\rho_0$ at
every step; it does not assume vanishing at unknown nearby points.
It constructs a possibly varying form, rather than assuming that
the deformation preserves the original fixed matrix.

\section{Application to all tested icosahedral points}

Use either standard complex three-dimensional rotation representation
of $A_5$. Multiplying by the parity sign gives the corresponding
rank-one reflection representation $\rho_0$ of $Q$ into $H_3$.
The checked deck implementer constructs it for each tested kernel.
The parity sign disappears on symmetric forms.

The key representation identity is
\[
              \Sym^2(W^*)\big|_G\cong\C[G/H_6].
\]
For completeness, put $\tau^2=\tau+1$. The character calculation is
\begin{center}
\begin{tabular}{lrrrrr}
\toprule
Class & $1$ & $2$ & $3$ & $5A$ & $5B$\\
\midrule
$\chi_W$ &3&$-1$&0&$\tau$&$1-\tau$\\
$\chi_{\Sym^2 W}$ &6&2&0&1&1\\
Six-point permutation character &6&2&0&1&1\\
\bottomrule
\end{tabular}
\end{center}
The second row follows from
$\chi_{\Sym^2 W}(g)=(\chi_W(g)^2+\chi_W(g^2))/2$; squaring swaps
the two five-cycle classes. The other rotation representation swaps
the last two entries in the first row, leaving the second row
unchanged. Complex representations of a finite group are semisimple
and determined by their characters, proving the identity.

Shapiro now gives
\[
 H^1(G,\Sym^2W^*)\cong H^1(H_6,\C)=0,
\]
using the finite cyclic homology established above. Restriction
to $G$ is injective in this first cohomology: the inflation kernel
is first cohomology of $Q/G=C_2$, which vanishes
in characteristic zero by averaging. Hence the hypothesis of
Theorem~\ref{thm:forms} holds.

\begin{corollary}\label{cor:peripheral}
For each tested standard icosahedral reflection point, every formal
complex deformation of its $Q$ representation preserves a
nondegenerate symmetric form. On the meridian's rank-one
nonidentity eigensummand, its commuting longitude has
$\lambda(t)^2=1$ and therefore constant eigenvalue $+1$ or $-1$.
\end{corollary}
\begin{proof}
The theorem constructs the form. Since $\mu^2=1$ and two is
invertible, $(I\pm\rho_t(\mu))/2$ are complementary idempotents.
Their images are free summands of ranks two and one, as at $t=0$.
Form invariance makes these summands orthogonal, with unit
restriction on the rank-one summand. The longitude preserves that
summand because it commutes with the meridian. Its scalar action
$\lambda(t)$ preserves the form there, so $\lambda(t)^2=1$.
The domain $\C[[t]]$ has only the two roots $\pm1$ of this equation.
\end{proof}

This strengthens the earlier pointwise orthogonal obstruction to a
formal-neighborhood obstruction under the stated meridian-square
condition. It does not assume the representation is constant or
that a common fixed point exists. It also does not treat knot-group
families with varying meridian eigenvalue. Those need not factor
through $Q$ and are precisely an option left open by the peripheral
construction in \cite[Sections 2.3--2.4]{gglw}. No formal family
or actual augmentation curve of that kind is constructed here.

\section{Replay, the generator probe and limits}

The cover producer uses conjugation on six explicitly listed
Sylow-five subgroups, with a frozen sixteen-cover cap, checkpointing
after each cover. The verifier instead constructs
$\operatorname{PSL}_2(\F_5)$ on $\mathbb P^1(\F_5)$. An explicit
group isomorphism and sheet bijection identify the coefficient
actions. A separate Fox-prefix matrix, with transported tree edges,
replays all Smith operations on the $84\times73$ cover matrices.
Direct augmentation cocycles reproduce each relevant field dimension
and the nonsplit characteristics. Orbifold relations and deck
implementers are checked separately. The PD parser and Smith replay
are shared. The receipt is \path{six_point_cover_verification.json}.

Known controls give index-six subgroup homology $\Z^7$ for $F_2$,
$\Z^3$ for $C_3*C_3$, and $\Z^7$ for $C_3*C_3*C_3$.
In characteristic five their Betti gains are respectively $5,3,7$.
Thus the two-generated free control attains the allowed maximum
five, while the known triple-trefoil control passes the new gate.

A different frozen probe tried to certify that the original cover
groups are two-generated. It substituted the old four-seed traces
into three-generator core presentations, then tested direct relators
and products of two signed relators with short conjugators under
96 explicit invertible free-group automorphisms. It tried 444
distinct candidates per target and found no primitive consequence.
Its verifier checks the source substitutions, inverse automorphisms
and positive/negative controls. The bounded misses are exploratory:
they prove neither rank two nor rank three, nor absence of a
primitive word in the full relation normal closure. The receipt
\path{core_generator_probe_verification.json} explicitly makes no
rank conclusion. Repeating longer word searches without a new
structural reason is not the next research gate.

The formal lifting argument is the uniform proof; finite matrices
verify its hypotheses on the retained maps. They are not formal or
external verification of that argument. No new MRC case is closed.
The full goal still requires a general proof or counterexample and
the outstanding literature audits, not completion of these controls
or exclusions.

\clearpage
\begin{thebibliography}{9}
\bibitem{gglw} Gao--Gu--Li--Wang, Meridional ranks of iterated
Whitehead doubles, preprint, 11 November 2026, Sections 2.3--2.4.
\url{https://arxiv.org/html/2611.15673v1}.
\bibitem{core} Silver--Traldi--Williams, Core groups (2024),
Section 7, including the Kelly--Wada double-cover normalization.
\url{https://arxiv.org/abs/2311.02048v2}.
\end{thebibliography}
\end{document}
